\documentclass[11pt,a4paper]{article}

\usepackage{fontspec}
\IfFontExistsTF{Times New Roman}{\setmainfont{Times New Roman}}{\setmainfont{Tinos}}
\IfFontExistsTF{Consolas}{\setmonofont{Consolas}[Scale=0.90]}{\setmonofont{Liberation Mono}[Scale=0.90]}

\usepackage[top=1.5cm, bottom=1.6cm, left=1.75cm, right=1.75cm]{geometry}
\setlength{\headheight}{18pt}

\usepackage{xcolor}
\usepackage{amsmath,amssymb}
\usepackage{tcolorbox}
\tcbuselibrary{skins}
\usepackage{booktabs}
\usepackage{tabularx}
\usepackage{listings}
\usepackage{fancyhdr}
\usepackage{lastpage}
\usepackage{enumitem}
\usepackage{titlesec}
\usepackage{hyperref}

\definecolor{primaryblue}{RGB}{18,65,115}
\definecolor{secondaryteal}{RGB}{14,102,110}
\definecolor{CodeBg}{RGB}{248,250,252}
\definecolor{DarkGray}{RGB}{55,55,55}
\definecolor{correctgreen}{RGB}{20,108,55}
\definecolor{rubricred}{RGB}{160,30,30}
\definecolor{boxbg}{RGB}{248,250,252}

\lstdefinestyle{pythonstyle}{
  language=Python,
  basicstyle=\ttfamily\scriptsize,
  keywordstyle=\color{primaryblue}\bfseries,
  commentstyle=\color{secondaryteal}\itshape,
  stringstyle=\color{rubricred},
  showstringspaces=false,
  columns=fullflexible,
  keepspaces=true,
  frame=single,
  rulecolor=\color{gray!40},
  backgroundcolor=\color{boxbg},
  breaklines=true,
  tabsize=4
}
\lstset{style=pythonstyle}

\hypersetup{
    colorlinks=true,
    linkcolor=primaryblue,
    urlcolor=secondaryteal,
    pdftitle={DSAI1003 - Objects and Classes Solutions and Marking Rubric},
    pdfauthor={Faculty of Data Science and Artificial Intelligence - National Economics University}
}

\pagestyle{fancy}
\fancyhf{}
\fancyhead[L]{\footnotesize\color{primaryblue}\textbf{DSAI1003: Fundamental Programming Concepts in Python}}
\fancyhead[R]{\footnotesize\color{correctgreen}\textbf{Official Solutions \& Marking Rubric: Chapter 9}}
\fancyfoot[L]{\scriptsize Faculty of Data Science and Artificial Intelligence -- National Economics University}
\fancyfoot[C]{\footnotesize Page \thepage\ of \pageref{LastPage}}
\fancyfoot[R]{\footnotesize\textbf{Total: 100 Points}}
\renewcommand{\headrulewidth}{0.5pt}
\renewcommand{\footrulewidth}{0.3pt}

\titleformat{\section}{\large\bfseries\color{primaryblue}}{\thesection}{1em}{}[\titlerule]
\titleformat{\subsection}{\normalsize\bfseries\color{secondaryteal}}{\thesubsection}{1em}{}
\titlespacing*{\section}{0pt}{3pt}{2pt}
\titlespacing*{\subsection}{0pt}{2pt}{1pt}

\newtcolorbox{solbox}[2][]{
  enhanced,
  colback=white,
  colframe=correctgreen,
  arc=1.5mm,
  boxrule=0.7pt,
  top=1.2mm,
  bottom=1.2mm,
  left=2.5mm,
  right=2.5mm,
  title=\textbf{\color{correctgreen}#2},
  coltitle=correctgreen,
  attach boxed title to top left={yshift=-2mm, xshift=2.5mm},
  boxed title style={colback=green!8, colframe=correctgreen, arc=1mm},
  #1
}

\newtcolorbox{rubricbox}[1][]{
  enhanced,
  colback=white,
  colframe=rubricred!80,
  arc=1.5mm,
  boxrule=0.7pt,
  top=1.2mm,
  bottom=1.2mm,
  left=2.5mm,
  right=2.5mm,
  title=\textbf{\color{rubricred}Marking Rubric \& Scoring Criteria},
  coltitle=rubricred,
  attach boxed title to top left={yshift=-2mm, xshift=2.5mm},
  boxed title style={colback=red!6, colframe=rubricred!80, arc=1mm},
  #1
}

\begin{document}

\begin{center}
  {\Large\textbf{NATIONAL ECONOMICS UNIVERSITY (NEU)}}\\[0.05cm]
  {\large\textbf{COLLEGE OF TECHNOLOGY --- FACULTY OF DATA SCIENCE \& AI}}\\[0.15cm]
  \rule{\textwidth}{1pt}\\[0.2cm]
  {\LARGE\bfseries\color{primaryblue}OFFICIAL SOLUTIONS \& MARKING RUBRIC}\\[0.1cm]
  {\large\textbf{DSAI1003 -- Fundamental Programming Concepts in Python}}\\[0.05cm]
  {\normalsize\textbf{Progress Exam: Objects and Classes (Chapter 9)} \quad | \quad Total: 100 Points}\\[0.15cm]
  \rule{\textwidth}{0.5pt}
\end{center}

\begin{tcolorbox}[colback=green!6, colframe=correctgreen, arc=1.5mm, top=1mm, bottom=1mm]
\small
\textbf{Note for Examiners and Teaching Assistants:}
This marking guide provides complete keys, rigorous explanations, step-by-step traces, reference Python implementations, and point-by-point deduction criteria to ensure fair and unified grading across all student cohorts.
\end{tcolorbox}

\vspace{0.1cm}
\section*{Section A: Object-Oriented Programming Foundations (20 points)}
\textit{2 points per question. Total = 20 points.}

\begin{center}
\renewcommand{\arraystretch}{1.25}
\small
\begin{tabularx}{\textwidth}{|c|c|X|c|}
\hline
\textbf{Q\#} & \textbf{Key} & \textbf{Conceptual Explanation \& Language Rule} & \textbf{Marks} \\
\hline
\textbf{Q1} & \textbf{B} & A class serves as an abstract blueprint or type definition defining attributes and behaviors, whereas an object is a concrete runtime instance created in memory from that blueprint. & 2 pts \\
\hline
\textbf{Q2} & \textbf{B} & In Python, \texttt{\_\_init\_\_} is the instance initialization method (constructor) automatically called after the object is allocated to set its initial state via \texttt{self}. & 2 pts \\
\hline
\textbf{Q3} & \textbf{B} & In Python method binding, calling \texttt{instance.method(arg)} translates behind the scenes to \texttt{Class.method(instance, arg)}. The explicit parameter \texttt{self} captures the receiving object instance. & 2 pts \\
\hline
\textbf{Q4} & \textbf{B} & Python follows the convention where a single leading underscore (e.g., \texttt{\_balance}) signals that an attribute is internal/private (encapsulation). It is not enforced by hardware or compiler barriers. & 2 pts \\
\hline
\textbf{Q5} & \textbf{C} & An accessor inspects or computes state without modifying the object (\texttt{getTotal()}), while a mutator alters the object's internal state variables (\texttt{addItem()}). & 2 pts \\
\hline
\textbf{Q6} & \textbf{B} & Variable assignment \texttt{c2 = c1} copies the object reference, not the object. Both names point to the same memory instance in the heap. Two clicks increment the shared count to 2. & 2 pts \\
\hline
\textbf{Q7} & \textbf{C} & The \texttt{is} operator tests identity (whether two references denote the identical memory address: \texttt{id(a) == id(b)}). Distinct constructor calls create distinct objects in heap, so \texttt{acc1 is acc2} is \texttt{False}. & 2 pts \\
\hline
\textbf{Q8} & \textbf{A} & Instance variables belong to specific object instances (defined using \texttt{self.\_var}), whereas class variables are declared in the class body and shared across all instances of that class. & 2 pts \\
\hline
\textbf{Q9} & \textbf{C} & Python automatically invokes \texttt{obj.\_\_str\_\_()} whenever \texttt{str(obj)} or \texttt{print(obj)} is called. It must return a human-readable informal string representation of the object. & 2 pts \\
\hline
\textbf{Q10} & \textbf{C} & In Python, \texttt{None} is an instance of \texttt{NoneType}. Attempting to call an arbitrary method on \texttt{None} causes Python to raise an \texttt{AttributeError: 'NoneType' object has no attribute '...'}. & 2 pts \\
\hline
\end{tabularx}
\end{center}

\newpage
\section*{Section B: Output Tracing, Object Mechanics \& Diagnostics (25 points)}

\subsection*{Q11. Exact Output of Multi-Object State Tracing (5 points)}
\begin{solbox}{Model Output for Q11}
\begin{verbatim}
c1: 8 | c2: 10
c1 is c3: True | c1 is c2: False
\end{verbatim}
\end{solbox}
\textbf{State Trace:}
\texttt{c1} starts with count 0, step 2; \texttt{c2} starts with count 0, step 5. \texttt{c3 = c1} creates an alias to \texttt{c1}. After \texttt{c1.click()}, \texttt{c1} count is 2. After \texttt{c2.click()}, \texttt{c2} count is 5. Calling \texttt{c3.boost(3)} scales \texttt{c1}'s step to 6. Calling \texttt{c3.click()} increments \texttt{c1} to $2 + 6 = 8$. Calling \texttt{c2.click()} increments \texttt{c2} to $5 + 5 = 10$. Finally, \texttt{c1 is c3} evaluates to \texttt{True} (shared object) while \texttt{c1 is c2} is \texttt{False}.\newline
\textbf{Scoring:} 2 pts for correct counter counts; 2 pts for identity booleans; 1 pt for exact formatting and labels.

\subsection*{Q12. Evaluate Class and Object Expressions (5 points --- 1 pt each)}
\begin{enumerate}[label=\alph*),itemsep=1.5pt]
\item \texttt{isinstance(it, Item)} evaluates to \textbf{\texttt{True}} (direct instance of class \texttt{Item}).
\item Exception class: \textbf{\texttt{AttributeError}} (attempting to read nonexistent attribute \texttt{price}).
\item \texttt{len(a)} = \textbf{\texttt{3}} (\texttt{b = a} aliases mutable list; mutating \texttt{b} updates \texttt{a} to \texttt{[10, 20, 30]}).
\item Value in \texttt{\_numerator}: \textbf{\texttt{1}} ($\gcd(4, 12) = 4 \implies 4 // 4 = 1$).
\item \texttt{bool(None)} evaluates to \textbf{\texttt{False}} (\texttt{None} evaluates to falsy in boolean context).
\end{enumerate}

\subsection*{Q13. Diagnose and Correct Three Object-Oriented Defects (5 points)}
\begin{center}
\renewcommand{\arraystretch}{1.22}
\small
\begin{tabularx}{\textwidth}{|p{4.5cm}|X|X|}
\hline
\textbf{Erroneous Code} & \textbf{Problem / Root Cause} & \textbf{Correction} \\
\hline
{\ttfamily\scriptsize class Account:}\newline{\ttfamily\scriptsize \ \ def deposit(amount):}\newline{\ttfamily\scriptsize \ \ \ \ self.\_balance += amount} & Signature lacks \texttt{self}. Calling \texttt{acc.deposit(100)} passes \texttt{acc} as 1st argument, causing \texttt{TypeError}. & Include \texttt{self} as 1st parameter:\newline{\ttfamily\scriptsize def deposit(self, amount):}\newline{\ttfamily\scriptsize \ \ \ \ self.\_balance += amount} (2 pts) \\
\hline
{\ttfamily\scriptsize class Counter:}\newline{\ttfamily\scriptsize \ \ def \_\_init\_\_(self): self.\_val = 0}\newline{\ttfamily\scriptsize \ \ def click(self): \_val += 1} & \texttt{\_val} lacks prefix \texttt{self.}; treated as unassigned local variable, raising \texttt{UnboundLocalError}. & Access attribute via \texttt{self}:\newline{\ttfamily\scriptsize def click(self):}\newline{\ttfamily\scriptsize \ \ \ \ self.\_val += 1} (2 pts) \\
\hline
{\ttfamily\scriptsize class Student:}\newline{\ttfamily\scriptsize \ \ def \_\_init\_\_(self, name, tags=[]):}\newline{\ttfamily\scriptsize \ \ \ \ self.\_tags = tags} & Dangerous mutable default \texttt{[]}, created once and shared across all instances without explicit list. & Use \texttt{None} default sentinel:\newline{\ttfamily\scriptsize def \_\_init\_\_(self, name, tags=None):}\newline{\ttfamily\scriptsize \ \ self.\_tags = [] if tags is None else list(tags)} (1 pt) \\
\hline
\end{tabularx}
\end{center}

\subsection*{Q14. Bank Account State Tracing with Overdraft \& Fee Rules (5 points)}
\begin{solbox}{Console Output for Q14}
\begin{verbatim}
Final Balance: 252.00
\end{verbatim}
\end{solbox}
\textbf{State Trace (3 pts):} Initial balance \$200.00 $\to$ \texttt{deposit(100.0)}: balance = \$300.00 $\to$ \texttt{withdraw(350.0)}: $350 > 300$ triggers \$10 penalty without withdrawal: balance = \$290.00 $\to$ \texttt{withdraw(50.0)}: $50 \le 290$, normal withdrawal: balance = \$240.00 $\to$ \texttt{add\_interest(5.0)}: $240 + 240 \times 0.05 = \mathbf{\$252.00}$. Output format: \texttt{252.00} (2 pts).

\subsection*{Q15. Object References vs Shallow Copies in Object Attributes (5 points)}
\begin{solbox}{Answers for Q15}
\begin{verbatim}
Cart 1: ['Book', 'Pen', 'Ruler', 'Eraser']
Cart 2: ['Book', 'Pen']
\end{verbatim}
\end{solbox}
\textbf{Technical Explanation (2 pts):} \texttt{c1} stored the reference to \texttt{raw} directly (\texttt{self.\_items = raw}). Any mutation on \texttt{raw} or \texttt{c1} modifies the shared list in place. \texttt{c2} was instantiated with \texttt{list(raw)}, which created an independent shallow copy containing \texttt{["Book", "Pen"]} at that moment; subsequent mutations to \texttt{raw} do not affect \texttt{c2}.

\newpage
\section*{Section C: Object-Oriented Design \& Modular Class Decomposition (25 points)}

\subsection*{Task 1: Public Methods, Parameters, Return Types, and Invariants (6 points --- 1.5 pts each)}
\begin{itemize}[itemsep=2pt]
\item \texttt{\_\_init\_\_(self, card\_id: str, initial\_balance: float = 0.0, daily\_cap: float = 8.0) -> None}:\\
  Initializes instance variables \texttt{\_card\_id}, \texttt{\_balance}, \texttt{\_daily\_cap}, \texttt{\_daily\_spent = 0.0}, and \texttt{\_rides = []}. Enforces invariants: $\text{\_balance} \ge 0$, $\text{\_daily\_cap} > 0$.
\item \texttt{top\_up(self, amount: float) -> float}:\\
  Precondition: \texttt{amount > 0}. Increases \texttt{\_balance} by \texttt{amount}. Returns new balance.
\item \texttt{tap\_ride(self, station: str, base\_fare: float) -> bool}:\\
  Precondition: \texttt{base\_fare > 0}. Computes effective fare subject to daily cap; deducts from balance if sufficient; logs ride and returns \texttt{True}; otherwise denies entry and returns \texttt{False}.
\item \texttt{reset\_daily(self) -> None}:\\
  Postcondition: Resets \texttt{\_daily\_spent} to \texttt{0.0} at start of a new calendar day.
\end{itemize}

\subsection*{Task 2: Implementation of Core Method \texttt{tap\_ride} (7 points)}
\begin{solbox}{Model Implementation for Task 2}
\begin{lstlisting}
def tap_ride(self, station, base_fare):
    # 1. Calculate remaining cap margin and effective fare
    remaining_cap = max(0.0, self._daily_cap - self._daily_spent)
    effective_fare = min(base_fare, remaining_cap)
    
    # 2. Guard against insufficient balance
    if self._balance < effective_fare:
        return False  # Turnstile gate entry denied
    
    # 3. Mutate internal state and record trip audit entry
    self._balance -= effective_fare
    self._daily_spent += effective_fare
    self._rides.append((station, effective_fare))
    return True
\end{lstlisting}
\end{solbox}
\textbf{Marking Rubric for Task 2:} 2 pts for capping formula; 2 pts for balance guard check returning \texttt{False}; 2 pts for updating balance, daily spent, and ride list; 1 pt for returning \texttt{True}.

\subsection*{Task 3: Hand-Trace a Transit Simulation Sequence (6 points --- 1.5 pts per step)}
Initial State: \texttt{balance = \$15.00}, \texttt{daily\_cap = \$8.00}, \texttt{daily\_spent = \$0.00}.
\begin{enumerate}[itemsep=1.5pt]
\item \texttt{tap\_ride("Central", 3.50)}:
  $\text{Fare} = \min(3.50, 8.00 - 0.00) = \mathbf{\$3.50} \implies \text{Balance} = 15.00 - 3.50 = \mathbf{\$11.50}, \quad \text{Daily} = \mathbf{\$3.50}$
\item \texttt{tap\_ride("WestEnd", 3.00)}:
  $\text{Fare} = \min(3.00, 8.00 - 3.50) = \mathbf{\$3.00} \implies \text{Balance} = 11.50 - 3.00 = \mathbf{\$8.50}, \quad \text{Daily} = \mathbf{\$6.50}$
\item \texttt{tap\_ride("NorthGate", 3.00)}:
  $\text{Remaining Cap} = 8.00 - 6.50 = 1.50 \implies \text{Fare} = \mathbf{\$1.50} \implies \text{Balance} = \mathbf{\$7.00}, \quad \text{Daily} = \mathbf{\$8.00}$ (Cap reached!)
\item \texttt{tap\_ride("Campus", 2.50)}:
  $\text{Remaining Cap} = 8.00 - 8.00 = 0.00 \implies \text{Fare} = \mathbf{\$0.00} \implies \text{Balance} = \mathbf{\$7.00}, \quad \text{Daily} = \mathbf{\$8.00}$ (Free ride)
\end{enumerate}
\textbf{Final State:} Final \texttt{\_balance} = \textbf{\$7.00}; Final \texttt{\_daily\_spent} = \textbf{\$8.00}; Total successful rides = \textbf{4}.

\subsection*{Task 4: Technical Analysis \& Object-Oriented Principles (6 points --- 3 pts each)}
\begin{enumerate}[label=\alph*),itemsep=1.5pt]
\item \textbf{Encapsulation and Invariant Enforcement (3 pts):}
  Encapsulation guarantees that internal variables (\texttt{\_balance}, \texttt{\_daily\_spent}) can only be altered through validated public methods. Direct external writes (e.g., \texttt{card.\_balance = -50.0}) would violate invariants, bypass the daily fare-capping logic, and prevent transaction logging.
\item \textbf{Pass-by-Object-Reference Semantics (3 pts):}
  Python passes objects by reference. When passing \texttt{card} to \texttt{batch\_process(card)}, the parameter receives a reference to the existing instance in heap memory, not a clone. Calling mutator methods inside the function alters the caller's original card object directly.
\end{enumerate}

\newpage
\section*{Section D: Applied Programming in Python (30 points)}

\subsection*{Problem 1: Commercial Bank Account with Tiered Overdraft (15 points)}

\begin{solbox}{Complete Reference Implementation for Problem 1}
\begin{lstlisting}
class CommercialAccount:
    """Represents a commercial checking account with overdraft facilities."""

    def __init__(self, acc_num, owner, balance=0.0, overdraft_limit=100.0):
        if balance < 0 or overdraft_limit < 0:
            raise ValueError("Balance and overdraft limit must be non-negative.")
        self._acc_num = str(acc_num)
        self._owner = str(owner)
        self._balance = float(balance)
        self._overdraft_limit = float(overdraft_limit)
        self._ledger = ["Account Opened"]

    def deposit(self, amount):
        if amount <= 0:
            raise ValueError("Deposit amount must be positive.")
        self._balance += amount
        self._ledger.append(f"Deposit: +${amount:.2f}")
        return self._balance

    def withdraw(self, amount):
        if amount <= 0:
            raise ValueError("Withdrawal amount must be positive.")
        if amount <= self._balance:
            self._balance -= amount
            self._ledger.append(f"Standard: -${amount:.2f}")
            return True
        elif amount <= self._balance + self._overdraft_limit:
            self._balance -= (amount + 10.0)  # $10 overdraft penalty fee
            self._ledger.append(f"Overdraft: -${amount:.2f} (Fee: $10.00)")
            return True
        else:
            self._ledger.append(f"Declined: -${amount:.2f}")
            return False

    def transfer_to(self, target_acc, amount):
        if self.withdraw(amount):
            target_acc.deposit(amount)
            return True
        return False

    def get_ledger(self):
        return list(self._ledger)

    def __str__(self):
        return f"Account[#{self._acc_num}, {self._owner}]: Balance = ${self._balance:.2f}"

def main():
    a1 = CommercialAccount("A101", "Alice Corp", 150.0, overdraft_limit=100.0)
    b2 = CommercialAccount("B202", "Bob Ltd", 50.0)
    a1.deposit(50.0)
    a1.withdraw(250.0)
    a1.transfer_to(b2, 30.0)
    print(a1)
    print(b2)
    for record in a1.get_ledger():
        print(f"  * {record}")

if __name__ == "__main__":
    main()
\end{lstlisting}
\end{solbox}

\begin{solbox}{Model Console Execution Output}
\begin{verbatim}
Account[#A101, Alice Corp]: Balance = $-100.00
Account[#B202, Bob Ltd]: Balance = $80.00
  * Account Opened
  * Deposit: +$50.00
  * Overdraft: -$250.00 (Fee: $10.00)
  * Overdraft: -$30.00 (Fee: $10.00)
\end{verbatim}
\end{solbox}

\begin{rubricbox}
\begin{itemize}[leftmargin=1.5em,itemsep=1.5pt]
\item \textbf{Constructor \& Attributes (3 pts):} Validation of non-negative balance and overdraft limit (1 pt); initializing attributes \texttt{\_acc\_num}, \texttt{\_owner}, \texttt{\_balance}, \texttt{\_overdraft\_limit}, and \texttt{\_ledger} with opening entry (2 pts).
\item \textbf{\texttt{deposit(self, amount)} (2 pts):} Guard check \texttt{amount > 0} (1 pt); balance update and formatted string logging (1 pt).
\item \textbf{\texttt{withdraw(self, amount)} (5 pts):} Guard check \texttt{amount > 0} (1 pt); standard branch deducting amount and logging (1.5 pts); overdraft branch deducting amount + \$10 fee within limit and logging (1.5 pts); declined branch logging without balance change (1 pt).
\item \textbf{\texttt{transfer\_to(self, target, amount)} (2 pts):} Atomic execution via withdrawal result triggering target deposit (2 pts).
\item \textbf{\texttt{\_\_str\_\_} \& \texttt{main()} (3 pts):} Correct formatting string representation matching spec (1 pt); driver test demonstration (2 pts).
\end{itemize}
\end{rubricbox}

\newpage
\subsection*{Problem 2: Immutable Rational Number Class with Operator Overloading (\texttt{Fraction}) (15 points)}

\begin{solbox}{Complete Reference Implementation for Problem 2}
\begin{lstlisting}
import math

class Fraction:
    """An immutable rational number represented by coprime numerator and denominator."""

    def __init__(self, numerator=0, denominator=1):
        if denominator == 0:
            raise ZeroDivisionError("Denominator cannot be zero.")
        if denominator < 0:
            numerator, denominator = -numerator, -denominator
        common = math.gcd(abs(numerator), denominator)
        self._num = numerator // common
        self._den = denominator // common

    def get_numerator(self): return self._num
    def get_denominator(self): return self._den

    def __str__(self):
        if self._num == 0: return "0"
        if self._den == 1: return f"{self._num}"
        return f"{self._num}/{self._den}"

    def __add__(self, other):
        n = self._num * other._den + other._num * self._den
        d = self._den * other._den
        return Fraction(n, d)

    def __sub__(self, other):
        n = self._num * other._den - other._num * self._den
        d = self._den * other._den
        return Fraction(n, d)

    def __mul__(self, other):
        return Fraction(self._num * other._num, self._den * other._den)

    def __eq__(self, other):
        if not isinstance(other, Fraction): return False
        return self._num == other._num and self._den == other._den

    def to_float(self): return self._num / self._den

def main():
    f1 = Fraction(1, 6); f2 = Fraction(1, 3)
    print(f"f1 = {f1}, f2 = {f2}")
    print(f"Addition: {f1} + {f2} = {f1 + f2}")          # 1/2
    print(f"Subtraction: {f1} - {f2} = {f1 - f2}")       # -1/6
    print(f"Multiplication: {f1} * {f2} = {f1 * f2}")    # 1/18
    print(f"Float value of {f1 + f2}: {(f1 + f2).to_float()}")

    f3 = Fraction(-2, -4); f4 = Fraction(3, -9)
    print(f"Reduced Fraction(-2, -4) = {f3} | Fraction(3, -9) = {f4}")
    print(f"Equality check (f1 + f2 == f3): {(f1 + f2) == f3}")

    try:
        Fraction(5, 0)
    except ZeroDivisionError as err:
        print(f"Caught expected error: {err}")

if __name__ == "__main__":
    main()
\end{lstlisting}
\end{solbox}

\begin{solbox}{Model Console Execution Output}
\begin{verbatim}
f1 = 1/6, f2 = 1/3
Addition: 1/6 + 1/3 = 1/2
Subtraction: 1/6 - 1/3 = -1/6
Multiplication: 1/6 * 1/3 = 1/18
Float value of 1/2: 0.5
Reduced Fraction(-2, -4) = 1/2 | Fraction(3, -9) = -1/3
Equality check (f1 + f2 == f3): True
Caught expected error: Denominator cannot be zero.
\end{verbatim}
\end{solbox}

\begin{rubricbox}
\begin{itemize}[leftmargin=1.5em,itemsep=1.5pt]
\item \textbf{Constructor Validation \& Invariants (4 pts):} Raising \texttt{ZeroDivisionError} when \texttt{denominator == 0} (1 pt); normalizing negative signs to numerator (1.5 pts); reducing to lowest terms via \texttt{math.gcd} and storing in \texttt{\_num} and \texttt{\_den} (1.5 pts).
\item \textbf{Immutability \& Accessors (1 pt):} Ensuring attributes are not mutated and providing \texttt{get\_numerator} / \texttt{get\_denominator}.
\item \textbf{String Representation \texttt{\_\_str\_\_} (2 pts):} Handling \texttt{0}, integer denominator \texttt{1}, and standard fraction format \texttt{f"\{num\}/\{den\}"}.
\item \textbf{Operator Overloading \texttt{\_\_add\_\_}, \texttt{\_\_sub\_\_}, \texttt{\_\_mul\_\_} (4 pts):} Correct mathematical formulas returning a fresh \texttt{Fraction} instance for each operation.
\item \textbf{Equality \texttt{\_\_eq\_\_} \& \texttt{to\_float()} (2 pts):} Comparing canonical reduced numerators and denominators (1 pt); float division (1 pt).
\item \textbf{\texttt{main()} Demonstration Driver (2 pts):} Testing addition, subtraction, multiplication, sign reduction, equality, and catching \texttt{ZeroDivisionError}.
\end{itemize}
\end{rubricbox}

\vfill
\begin{center}\rule{0.5\textwidth}{0.4pt}\\[1mm]
\textbf{\color{primaryblue}--- END OF SOLUTIONS \& RUBRIC ---}
\end{center}

\end{document}
